\documentclass[../../main.tex]{subfiles}
\begin{document}


{\bf [解]}

$e(\theta) = \cos\theta + i\sin\theta$ とおく．3数 $\alpha_1, \alpha_2, \alpha_3$ に対して
\begin{align} 
  \alpha_k = r_k e(\theta_k) \quad (r_k, \theta_k \in \mathbb{R}) 
\end{align}
とする．ただし$k=1, 2, 3$とする．$\alpha_k$のなかに$0$があるかどうかで場合分けして考える．

\subparagraph{$\alpha_k=0$なる$k$がある時}

対称性から $\alpha_1=0$ として考える．題意より$\alpha_1 \ne \alpha_2, \, \alpha_1 \ne \alpha_3$ だから $r_2, r_3 > 0$ である．
題意の条件から
\begin{align} 
  \alpha_2 \alpha_3 = \alpha_1, \alpha_2, \alpha_3 
\end{align}
だが，左辺は非零なので$\alpha_1$ではあり得ず
\begin{align}
 \alpha_2 \alpha_3 = \alpha_2, \alpha_3 \\
\therefore
 \alpha_2 = 1 \text{ or } \alpha_3 = 1
\end{align}
となる．対称性から$\alpha_2=1$の時を考えよう．
さらに題意の条件から
\begin{align} 
  \alpha_3^2 = 0, 1, \alpha_3 
\end{align}
である．$\alpha_k$が互いに異なるので$\alpha_3=-1$が解である．以上から
\begin{align}
 (\alpha_1, \alpha_2, \alpha_3) = (0, 1, -1)
\end{align}
であることが必要である．
この3数は題意の条件を満たして十分である．

\subparagraph{全ての$\alpha_k$が非零の場合}

この時 $r_k > 0$ であり, 題意から $r_k = 1$ が必要である．まずこれを示そう．
題意より$r_3^2$が$r_1,r_2,r_3$のいずれかに等しいから，これを繰り返し用いて任意の自然数$n$に対して$r_3^n$も$r_1,r_2,r_3$のいずれかに等しい．ところが$r_3\neq 0$より$r_3 \neq 1$の時は$r_3$は$0$に収束または無限大に発散してしまい$r_1,r_2,r_3$と同じにならない．従って$r_3=1$が必要である．$r_1, r_2$についても同様だから$r_k=1$である．

この元で題意の条件
\begin{align}
 \alpha_1 \alpha_2 = \alpha_1, \alpha_2, \alpha_3 \label{eq:2}
\end{align}
を満たす$\alpha_k$を考えよう．

$\alpha_1 \alpha_2 = \alpha_1$ or $\alpha_1 \alpha_2 = \alpha_2$ の時は $\alpha_1 = 1$ or $\alpha_2 = 1$ となる．
対称性から $\alpha_1 = 1$ とする．
この時, $\alpha_2 \ne 1, \, \alpha_3 \ne 1, \, \alpha_2 \ne \alpha_3$ より, 題意のその他の条件は
\begin{align} 
  \alpha_2 \alpha_3 = 1, \quad \alpha_2^2 = \alpha_3 
\end{align}
となり, 
\begin{align}
 \{ \alpha_2, \alpha_3 \} = \left\{ e\left(\dfrac{2}{3}\pi\right), e\left(\dfrac{4}{3}\pi\right) \right\}
\end{align}
を得る．逆にこの時題意の条件は全て満たして十分である．すなわち
\begin{align}
 \left(\alpha_1,\alpha_2,\alpha_3\right) = \left(1, e\left(\dfrac{2}{3}\pi\right), e\left(\dfrac{4}{3}\pi\right) \right) \label{eq:3}
\end{align}
は解の一つである．

最後に, \cref{eq:2} において $\alpha_1 \alpha_2 = \alpha_3$ の時, \cref{eq:3}の場合をのぞいて, $\alpha_1 \ne 1, \alpha_2 \ne 1, \alpha_3 \ne 1$ として良い．
この時同様に
\begin{align}
\begin{cases}
\alpha_2 \alpha_3 = \alpha_1 \\
\alpha_3 \alpha_1 = \alpha_2
\end{cases} \label{eq:4}
\end{align}
となるしかないから，辺々かけて, 
\begin{align}
 \left(\alpha_1\alpha_2\alpha_3\right)^2 =  \left(\alpha_1\alpha_2\alpha_3\right)
\end{align}
である．両辺 $\alpha_1 \alpha_2 \alpha_3\, (\neq 0)$ でわって,
\begin{align} 
 \alpha_1 \alpha_2 \alpha_3 = 1 
\end{align}
となる．\cref{eq:4}を代入して
\begin{align}
 \alpha_1^2 = \alpha_2^2 = \alpha_3^2 = 1 \\
\therefore
 \alpha_k = \pm 1
\end{align}
だが，これは仮定（$\alpha_k \neq 1$および$\alpha_k$が互いに異なる）に反するので矛盾．

以上よりもとめる3数の組み合わせは
\begin{align} 
(0, 1, -1), \quad \left(1, -\dfrac{1}{2}+\dfrac{\sqrt{3}}{2}i, -\dfrac{1}{2}-\dfrac{\sqrt{3}}{2}i\right)  
\end{align}
の二つである．

\end{document}
